\chapterimage{onboarding.pdf}
\chapterimagecaption{\textit{The School of Athens} (1511) by Raphael}
\chapter{Onboarding}
    \section{Onboarding Checklist} \index{Onboarding}
Welcome to McGill! Below is a checklist for new students. Make sure to go over these items! This will make your life much easier.

\begin{itemize}[
    label={$\square$},
    leftmargin=1.8em,
    itemsep=1em
]

    \item Register as a \textbf{``full-time''} student by adding the Registration Confirmation course (REGN RCGR) for both the Fall and Winter terms. The registration period is until \textbf{August 14th, 2026}. Classes begin on \textbf{August 31st}. The Add/Drop deadline is \textbf{September 15th}.
    
    \item MSc students register for 30 research credits towards their thesis and 15 credits of courses. Typically, \textbf{PHYS 693} is taken in the 1st year with courses; \textbf{PHYS 690D1} and \textbf{PHYS 690D2} are taken in the 2nd year.

    \item PhDs take two classes at the 600+ level. Students who complete their MSc at McGill can carry over one of the 600+ level courses. PhDs must also pass the \textbf{PhD Preliminary Exam}.

    \item Additionally, new students must also take two professional development courses, \textbf{PHYS 601 \& 602}. 

    \item Check your official McGill e-mail account (first.last@mail.mcgill.ca) for announcements, scholarship opportunities and formal paperwork.

    \item Request a \textbf{Tuition Deferral} in Minerva to avoid interest charges on your student account. Tuition supplements can only be paid out on September 1st.

    \textit{Minerva > My Financial Aid \& Awards > Defer Payment of Tuition \& Fees.}

    \item \textbf{Conditions of Admission:} Submit your final official transcripts/proof of graduation to Enrolment Services. Many universities are now providing electronic copies. They can be sent to \href{mailto:officialschooldocs@mcgill.ca}{officialschooldocs@mcgill.ca}.

    \item \textbf{International students:}  Keep the department updated and let them know when you receive your \textbf{CAQ} and \textbf{Study Permit}. Use \href{https://www.mcgill.ca/legaldocuments/how}{this} site for information on uploading your documents. When you arrive in Canada, you will open a Canadian bank account and update your mailing address in Minerva. You have options for receiving your research stipend before you arrive in Canada. \\    

    \textbf{\color{ocre}Note:} When you arrive in Canada, you can request an adjustment (or refund) for the \textbf{International Health Insurance premium}, for the months that you were not in Canada. \index{Health Insurance!International Health Insurance}

    \item \textbf{Canadian students:} Update your personal banking information in Minerva to set up your \textbf{Direct Deposit}. 
    
    \item Send your signed \textbf{TA contract} and personal data form to Kim (\href{mailto:ao.physics@mcgill.ca}{ao.physics@mcgill.ca}), so that she can create your TA appointment in Workday.

     \item \textbf{Finding a desk:} You must make sure to set up an office space from where you can work, this procedure may vary slightly from student to student depending on their group, see the next section for more details. Email Kim when you have a desk, for access.
     
     \item \textbf{Building access:} As an active member of the physics community, you should make sure you have keycard access to your lab, office, and the Rutherford Physics Building. Since Louise has retired, reach out to Robert Turner (\href{mailto:robert.turner@mcgill.ca}{robert.turner@mcgill.ca}) to request keycard access.

    \item Email the Graduate Coordinator, Sam Maroussis (\href{mailto:Graduate.Physics@McGill.ca}{Graduate.Physics@McGill.ca}) if you have any questions!

\end{itemize}

\medskip

\textbf{\color{ocre}Note:} you will also receive this checklist from Sam to help you get set up.

\subsection*{McGill E-Mails}\index{Emails}
\begin{itemize}
    \item You have a McGill email: 
    \href{mailto:first.last@mail.mcgill.ca}{first.last@mail.mcgill.ca}.

    \item You can get an account on the department Linux network (ask your supervisor) which will give you an @physics.mcgill.ca e-mail.

    \item If you are a teacher assistant (TA), you will also be assigned an e-mail which is 
    \href{mailto:@mcgill.ca}{@mcgill.ca}. Note that this is temporary! It will go away if you stop being a TA, such as during the summer or during terms where you have a TA buyout from your supervisor. 
    \item To avoid the failure mode where you cannot access emails sent to your @mcgill.ca emails during times when this email is deactivated, consider setting up Outlook forwarding from your @mcgill.ca to @mail.mcgill.ca account.
\end{itemize}

It’s really important to check your McGill email (@mail.mcgill.ca) -- this is where announcements about things like funding, workshops will come. You should also use it for communication within the university (you may not get a response otherwise). Department announcements can sometimes go to your physics email address \href{mailto:grads@physics.mcgill.ca}{grads@physics.mcgill.ca} (if you have one). You can have this redirected to your McGill email address by adding a \texttt{.forward} file in your home directory, or if you don’t have a physics account you can email Juan Gallego \href{mailto:operator@physics.mcgill.ca}{(operator@physics.mcgill.ca)} to make one.

\subsection*{ID Cards}\index{ID Cards}
To get your McGill ID, follow the steps on \href{https://www.mcgill.ca/student-records/personal-information/id}{this} website. You just need to upload your photo using the upload tool, wait for it to get approved, and then collect it from Service Point.

\subsection*{Building Access}\index{Building Access}
There are many different types of access within the Rutherford Physics Building. Access can be requested only after the ID card has been physically picked up by the student. They must send Robert Turner an email indicating their name, email, and the room for which access is requested, while CCing their supervisor. 

All graduate students are guaranteed after-hour access to the building. Send another mail to Robert (without CCing your supervisor) with your McGill ID and a PIN. To open Rutherford doors after-hours, tap your keycard on the reader in the main entrance and type the \textbf{\color{ocre} last four digits} of your student ID number or the PIN of your choosing. Access to specific offices must be requested directly from the administration. 

In case your work requires access to your group's laboratory space, contact your supervisor and the administration. To access specific research facilities and creative spaces within the Rutherford and Wong buildings, such as the Wet Lab or the Makerspace, you must contact the technician responsible for the site to obtain access. Rooftop access is restricted to individuals working with the Anna McPherson observatory or who work with research materials stored on the roof.


\subsection*{Getting a Desk} \index{Finding a Desk}
The most direct and most commonly used way to secure a desk in the Rutherford Physics Building offices is to pick the office of your choice and check with the current occupiers if there are any desks available. Once you have found a free desk, make sure to claim it by attaching a name tag to it: simply write your name and contact information. Afterwards, contact Kim Metera (\href{mailto:ao.physics@mcgill.ca}{ao.physics@mcgill.ca}) with your room and desk number.

\begin{plainexercise}
 Note that even if you claim a desk by adding your name tag, you must make sure to actively use it otherwise it is entirely possible someone else will claim it as their own.    
\end{plainexercise}
\vspace{-5pt}
Some research groups have desks inside their laboratories, in which case you should check with your group mates to check if any of those desks are available. Do \textbf{NOT} claim multiple desks in the building, desk space is limited and cases of double-desking will inevitably lead you to get kicked out of one of your desks.

Recently, the MGAPS council is mounting a Graduate Study Room in ERB 433, featuring communal desks equipped with monitors. This initiative is meant to accomodate any students that do not have a permanent desk and are on the desk waitlist, which is especially the case for incoming students.

\subsection*{Room Booking} \index{Room Booking}
To book a room in the Rutherford, contact Maha (\href{mailto:chairsec.physics@mcgill.ca}{chairsec.physics@mcgill.ca}) for now. All room bookings are detailed on \href{https://calendar.google.com/calendar/u/0/embed?height=600&wkst=1&bgcolor=%23ffffff&ctz=America/Toronto&src=bG91aXNlLmRlY2VsbGVzQG1jZ2lsbC5jYQ&src=cGh5c2ljc3Jvb21zQGdtYWlsLmNvbQ&color=%23039BE5&color=%23039BE5}{this calendar}, consult it to verify which rooms are available. Note that if you wish to book the Rutherford Auditorium (ERB 112), you must submit a request to the McGill Faculty of Science.

\subsection*{Safety and Training}\index{Safety}
The Department Safety Committee (DSC) is responsible for ensuring safety standards within the Rutherford Physics Building. Primarily, this includes laboratory safety, disposal of laboratory materials and junk, construction, and building operation (e.g. shut downs and leaks). The Chair and Safety Officer, Michael Hilke, is primarily responsible for ensuring these safety standards. Matters such as junk removal, on-going construction, and maintenance are usually handled by the Building Director, Andreas Warburton. 

The Environmental Health and Safety (EHS) is the entity responsible for training and certifying laboratory personnel. Any research involving laboratory work requires EHS certification and safety training, and some types of work (e.g. radiation, lasers) may require further training and permits.

For more on the role of the DSC, please see \href{https://www.mcgill.ca/ehs/policies-and-safety-committees/health-and-safety-committees/departmental-safety-committees}
{here}. For more information regarding the EHS, see \href{https://www.mcgill.ca/ehs/training}{here}. Also, contact your supervisor for more information regarding your own laboratory safety procedures and certifications.

\section{Relevant Contacts} \label{contacts}

\subsection*{Administration, Advisors, and Committee Chairs}
Below is a comprehensive list of the administration team, their roles, and contacts (as of August 2026).

\vspace{-5 pt}
    \begin{table}[h!]
        \hspace*{-1.5cm}
    \centering
    \renewcommand{\arraystretch}{1.2}
    \arrayrulecolor{ocre}
    \begin{tabular}{|>{\columncolor{ocre!20}}lc>{\columncolor{ocre!20}}cc|}
    \hline
    \rowcolor{ocre!60}
  {\color{white}\large\sffamily \textbf{Role}} & {\color{white}\large\sffamily \textbf{Name}}  & {\color{white}\large\sffamily \textbf{Office}}  & {\color{white}\large\sffamily \textbf{Email Address}} \\
        Department Chair & Ken Ragan & ERB 108 & chair.physics@mcgill.ca\\
        Building Director & Andreas Warburton & ERB 343 & buildingdirector.physics@mcgill.ca \\
        Safety Officer & Michael Hilke & ERB 408 & michael.hilke@mcgill.ca \\
        Administrative Officer & Kim Metera & ERB 108 & ao.physics@mcgill.ca \\
        Graduate Studies Coordinator & Sam Maroussis & ERB 108 & graduate.physics@mcgill.ca \\
        Graduate Studies Advisor & Nikolas Provatas & ERB 218 & nikolaos.provatas@mcgill.ca \\
        Undergraduate Advisor & Maha Ibrahim & ERB 223 & advising.physics@mcgill.ca \\
        Reimbursements \& Financials & Sonal Patel & ERB 443 & scipod1.managers@mcgill.ca \\
        Administrative Coordinator & Eddie Del Campo & ERB 108 & eduardo.delcampo@mail.mcgill.ca \\
        Administrative Coordinator & Alba Furfaro & ERB 108 & alba.furfaro@mcgill.ca \\
        Computer System Operator & Juan Gallego & ERB 331 & juan.gallego@mcgill.ca \\
        Library Liaison & April Colosimo & & april.colosimo@mcgill.ca \\
        EDI Committee Chair & Jason Hessels & ERB 215 & jason.hessels@mcgill.ca \\
        Inreach/Outreach Coordinator & Eesha Das Gupta & TSI 030A & eesha.dasgupta@mcgill.ca \\
        \hline
\end{tabular}
\arrayrulecolor{black}
\end{table}

\subsection*{Technicians and Professionals}   
Below is a table containing all technicians and research professionals affiliated with the Physics Department (as of August 2026). In case you require assistance, contact your supervisor to determine the best way to proceed and the best suited research professional.

\vspace{-5 pt}
 \begin{table}[h!]
        \hspace*{-1.7cm}
    \centering
    \renewcommand{\arraystretch}{1.2}
    \arrayrulecolor{ocre}
    \begin{tabular}{|>{\columncolor{ocre!20}}lc>{\columncolor{ocre!20}}cc|}
    \hline
    \rowcolor{ocre!60}
  {\color{white}\large\sffamily \textbf{Role}} & {\color{white}\large\sffamily \textbf{Name}}  & {\color{white}\large\sffamily \textbf{Office}}  & {\color{white}\large\sffamily \textbf{Email Address}} \\

        Linux Support and Physics Servers & Juan Gallego & ERB 331 & juan.gallego@mcgill.ca \\
        Machine Shop Supervisor & Pascal Bourseguin & ERB 004 & pascal.bourseguin@mcgill.ca \\
        Machine Shop Technician & Marc Giroux & ERB 004 & marc.giroux@mcgill.ca \\
        Electronics Design & Eamon Egan & ERB 228 & eamon.egan@mcgill.ca \\
        Rutherford Sample Preparation Lab &  Robert Gagnon & ERB 421A & robert.gagnon@mcgill.ca \\
        Technician: Vacuum Systems & John Smeros & ERB 420B & john.smeros@mcgill.ca \\
        Technician: Electronics Support & Richard Talbot & ERB 422A & richard.talbot@mcgill.ca \\
        Technician: Windows Support & Janney Wu & WON 0130 & janney.wu@mcgill.ca \\
        Undergraduate Lab Manager & Robert Turner & WON 0121 & robert.turner@mcgill.ca \\
        Technician: Undergrad Teaching Lab  & Anwesha Basu & Trottier 3080 & anwesha.basu@mcgill.ca \\
        Technician: Undergrad Teaching Lab  & Robert Idsinga & WON 0101 & robert.idsinga@mcgill.ca \\
        Technician: Undergrad Teaching Lab  & Philippe Movafegh & WON 0130 & philippe.movafegh@mcgill.ca \\
        Technician: Undergrad Teaching Lab & Brandon Ruffolo & WON 0130 & brandon.ruffolo@mcgill.ca \\
        Nanotools Manager & Zhao Lu & ERB 014 & zhao.lu@mcgill.ca \\
        Nanotools Equipment Technologist & John Li & ERB 013 & john.li@mcgill.ca \\
        Nanotools Research Professional & Peng Yang & ERB 013 & peng.yang2@mcgill.ca \\
        \hline
\end{tabular}
\arrayrulecolor{black}
\end{table}
\newpage

\subsection*{Wider University} 
Below is a list of contacts for individuals within McGill that could prove useful (as of August 2026).
\vspace{-5 pt}
 \begin{table}[h!]
        \hspace*{-1.7cm}
    \centering
    \renewcommand{\arraystretch}{1.2}
    \arrayrulecolor{ocre}
    \begin{tabular}{|>{\columncolor{ocre!20}}lc>{\columncolor{ocre!20}}cc|}
    \hline
    \rowcolor{ocre!60}
  {\color{white}\large\sffamily \textbf{Role}} & {\color{white}\large\sffamily \textbf{Name}}  & {\color{white}\large\sffamily \textbf{Site}}  & {\color{white}\large\sffamily \textbf{For Appointments}} \\

        Dean of Students & Tony Mittermaier & \href{https://mcgill.ca/deanofstudents/}{Link}& deanofstudents@mcgill.ca \\
        Dean of Graduate and Post-Doctoral Studies & Marie-Hélène Pennestri & \href{https://mcgill.ca/gps/}{Link}& Deanadm.gps@mcgill.ca \\
        Dean of Science & Alana Watt & \href{https://mcgill.ca/science/}{Link} & secdean.science@mcgill.ca \\
        \hline
\end{tabular}
\arrayrulecolor{black}
\end{table}


\section{Getting Involved}
\subsection{McGill Graduate Association of Physics Students (MGAPS)} \index{MGAPS}
The McGill Graduate Association of Physics Students (MGAPS) is the official division of all graduate physics students at McGill University. In other words, every graduate student is an MGAPS member. A \$10 fee for MGAPS membership is automatically billed into the graduate student tuition by the Post-Graduate Student Society (comprised of all graduate student associations in McGill University) into the MGAPS yearly budget for the organization of actions and events that aim to improve graduate students' experiences. The MGAPS council is a group of elected volunteers that serves as the governing body of MGAPS and is comprised of a president, vice-presidents, officers, representatives, and liaisons (for a full list of these roles and their descriptions see the \href{https://mgaps.physics.mcgill.ca/people/}{list of current council members} or the \href{https://mgaps.physics.mcgill.ca/files/MGAPS_Constitution_2025-01-28.pdf}{constitution}). The MGAPS council has as primary objectives the maintenance and improvement of wages, infrastructure, and enrichment for MGAPS members. 

\begin{plaincorollary}
    Becoming involved with the MGAPS council and MGAPS events is one of the most direct ways for graduate physics students to meet new people, develop leadership skills, and learn how to navigate new research environments.
\end{plaincorollary}
To become involved with MGAPS and the council, the council meets on a \textbf{biweekly basis} and organizes a series of weekly and recurring events, all details of which are noted in the \href{https://mgaps.physics.mcgill.ca/community/events/}{MGAPS calendar}). All MGAPS council activities and meetings are open to all MGAPS members. The MGAPS council also maintains the \textbf{Graduate Student Lounge (ERB 322)}, located in the 3rd floor of the Rutherford Physics Building, as a social space for all graduate students and post-docs in physics, and the lounge operates as the primary locations for council meetings and small MGAPS events. Any member of MGAPS can assume a position within the MGAPS council by winning the elections hosted in the Fall MGAPS General Assemblies, and any vacant positions will be again open for election in the Winter MGAPS General Assemblies as per the \href{https://mgaps.physics.mcgill.ca/files/MGAPS_Constitution_2025-01-28.pdf}{constitution}.
\begin{plainexercise}
    The MGAPS council is involved with various bodies both in and out of McGill University, including: Physics department administration, the Physics research centers (CHEP, CPM, and TSI), the Equity, Diversity, and Inclusion Committee, the Outreach Committee, the Teaching Support Union at McGill, the Post-Graduate Student Society, the Faculty of Science, and industry.
\end{plainexercise}
So, if you are interested, get involved with the MGAPS council! In order to contact any MGAPS council see the contacts listed below or see the \href{https://mgaps.physics.mcgill.ca/people/}{list of current council members} for specific contacts. Join the MGAPS Slack \href{https://mcgillphysics-rag1431.slack.com/join/shared_invite/zt-3vlrg6luf-k0KqSG3LdCQ2V2P96XF_kQ#/shared-invite/email}{here}.

\vspace{-5 pt}
 \begin{table}[h!]
    \centering
    \renewcommand{\arraystretch}{1.2}
    \arrayrulecolor{ocre}
    \begin{tabular}{|l>{\columncolor{ocre!20}}cc|}
    \hline
    \rowcolor{ocre!60}
  {\color{white}\large\sffamily \textbf{Role}} & {\color{white}\large\sffamily \textbf{Name}} &{\color{white}\large\sffamily \textbf{Email}} \\

        President and Vice Presidents & MGAPS & mgaps.exec@physics.mcgill.ca \\
        \hline 
\end{tabular}
\arrayrulecolor{black}
\end{table}
Note that it may be more practical to reach out to specific contacts.

\subsection{McGill Society of Physics Students (MSPS)}
The McGill Society of Physics Students (MSPS) is the official division of all undergraduate physics students at McGill University. Therefore, all undergraduate physics students are members of MSPS, and MSPS is an organization within the Students’ Society of McGill University. The MSPS council is a volunteer group of undergraduates that is elected as the representative body of MSPS, and a full list of the roles and their descriptions can be found in the \href{https://www.msps.physics.mcgill.ca/}{MSPS website}. The MSPS council maintains the Undergraduate Student Lounge as a social space for undergraduates, and MSPS council meetings take place in this lounge. The MSPS collaborates with MGAPS, the Equity, Diversity, and Inclusion Committee, and the Outreach Committee in the organization of events. The MSPS council is also the primary body within the physics department that is responsible for ordering McGill Physics merchandise for both undergraduate and graduate students, orders are handled seasonally by MSPS and can be ordered through the \href{https://www.msps.physics.mcgill.ca/}{MSPS website}. To contact the MSPS council, see the contacts listed below or see the \href{https://www.msps.physics.mcgill.ca/}{MSPS website} for specific contacts.

\vspace{-5 pt}
 \begin{table}[h!]
    \centering
    \renewcommand{\arraystretch}{1.2}
    \arrayrulecolor{ocre}
    \begin{tabular}{|l>{\columncolor{ocre!20}}cc|}
    \hline
    \rowcolor{ocre!60}
  {\color{white}\large\sffamily \textbf{Role}} & {\color{white}\large\sffamily \textbf{Name}}   & {\color{white}\large\sffamily \textbf{Contact}} \\

        President and Vice Presidents & MSPS & msps.physics@mcgill.ca \\
        \hline
\end{tabular}
\arrayrulecolor{black}
\end{table}
Note that it may be more practical to reach out to specific contacts.

\subsection{Equity, Diversity, and Inclusion (EDI) Committee} \index{Coordinators!EDI}
The Equity, Diversity, and Inclusion (EDI) Committee in the Department of Physics at McGill is a group of students, postdoctoral fellows, staff, and faculty members that aims to create a safe and inclusive working environment. The EDI Committee seeks to achieve this by organizing social and educational events, participating in outreach with CÉGEPs, maintaining EDI-related resources, and organizing with the physics administration team to improve departmental policies and conditions. The committee is led by a chair and composed of members, all of which are listed in the \href{https://www.physics.mcgill.ca/edi/}{EDI Committee website}, and frequently collaborate with MGAPS and MSPS. 

To become involved with committee activities, students can attend the \textbf{weekly Physics and Society Discussion Group} to discuss topics suggested by students and special EDI events announced by the committee. The committee also organizes a \textbf{weekly lunch for women and non-binary people} in physics in the Graduate Student Lounge (ERB 322), and the lunch is catered on a biweekly basis. Finally, the EDI Committee organizes an \href{https://www.physics.mcgill.ca/edi/library.html}{Equity Library} on the 2nd floor of the Rutherford Physics Building that covers equity issues in Canada and the United States. See the MGAPS Newsletter or the \href{https://www.physics.mcgill.ca/edi/events.html}{EDI Committee events page} for more details on recurring EDI events and resources.

To become a member of the EDI Committee, students may apply for one of the two paid \textbf{Graduate EDI Coordinator} positions open for graduate students (the pay is equivalent to a teaching assitant position) or be elected for the EDI Officer position within the MGAPS council. To contact the EDI Committee, see the contacts listed below or see the \href{https://www.physics.mcgill.ca/edi/}{EDI Committee website} for more specific contacts.

\vspace{-5 pt}
 \begin{table}[h!]
        \hspace*{-1.2cm}
    \centering
    \renewcommand{\arraystretch}{1.2}
    \arrayrulecolor{ocre}
    \begin{tabular}{|>{\columncolor{ocre!20}}lc>{\columncolor{ocre!20}}c|}
    \hline
    \rowcolor{ocre!60}
  {\color{white}\large\sffamily \textbf{Role}} & {\color{white}\large\sffamily \textbf{Name}}   & {\color{white}\large\sffamily \textbf{Contact}} \\

        EDI Committee Chair & Jason Hessels & jason.hessels@mcgill.ca \\
        Graduate EDI Coordinator & Angela Sabzevari Gonzalez & angela.sabzevarigonzalez@mail.mcgill.ca \\
        Graduate EDI Coordinator & Maryam Rahbaran & maryam.rahbaran@mail.mcgill.ca \\
        EDI Officer & Jack Faller & jacob.faller@mail.mcgill.ca \\
        \hline
\end{tabular}
\arrayrulecolor{black}
\end{table}

\subsection{Outreach Committee: McGill Physics and TSI Outreach} \index{Coordinators!Outreach}
The McGill Physics and TSI Outreach Committee is the primary organizing body behind outreach and science communication initiatives in the Physics department, and is composed of the McGill Physics and the Trottier Space Institute (TSI) outreach teams. The Outreach Committee aims to communicate Physics and Astronomy research to the general public, students, and the McGill community, and counts on the participation of paid outreach coordinators (the salary being roughly equivalent to a teaching assistantship) and numerous volunteers. 

To become members of the Outreach Committee, graduate students can become \textbf{outreach coordinators} by applying for the positions. Alternatively, graduate students can also become members by being elected as VP Professional Development in the MGAPS Council, who acts as a liaison between the committee and the MGAPS council. The Outreach Committee frequently calls upon volunteers from both graduate and undergraduate student bodies to organize its events, and it is an excellent way to begin involvement in outreach. To contact the Outreach Committee, see the contacts listed below or see the \href{https://physics-tsi-outreach.physics.mcgill.ca/index.html#team}{Outreach Committee website} for more specific contacts. Join the Outreach Slack to check for department-wide activity.

\vspace{-5 pt}
 \begin{table}[h!]
        \hspace*{-2.7cm}
    \centering
    \renewcommand{\arraystretch}{1.2}
    \arrayrulecolor{ocre}
    \begin{tabular}{|>{\columncolor{ocre!20}}lc>{\columncolor{ocre!20}}c|}
    \hline
    \rowcolor{ocre!60}
  {\color{white}\large\sffamily \textbf{Role}} & {\color{white}\large\sffamily \textbf{Name}}   & {\color{white}\large\sffamily \textbf{Contact}} \\

        Outreach Committee Chair & Thomas Brunner & outreach@physics.mcgill.ca \\
        Outreach Coordinator & Eesha Das Gupta & eesha.dasgupta@mcgill.ca\\
        Public Talks and Social Media Coordinator & Amélie Chiasson David & amelie.chiassondavid@mail.mcgill.ca\\
        Public Observing Nights & Anna I. McPherson Observatory & observatory@physics.mcgill.ca  \\
        Hackathon Coordinator &  Brendan Barry & hackathon@physics.mcgill.ca \\
        Space Explorers Coordinator & Patrick Horlaville & patrick.horlaville@mail.mcgill.ca \\
        Science in Space Coordinator & Georgia Mraz & georgia.mraz@mail.mcgill.ca \\
        MGAPS Liaison & Catherine Boisvert & catherine.boisvert4@mail.mcgill.ca\\
        \hline
\end{tabular}
\arrayrulecolor{black}
\end{table}

\newpage
\subsubsection{Public Talks} \index{Coordinators!Public Talks}
The Physics Department offers public talks roughly monthly during the academic year. Speakers are often professors, permanent researchers, postdocs, or senior PhD students-including researchers from local institutions and visiting scholars alike. The McGill Physics and TSI Public Talks are organized by a paid Outreach Coordinator (see contact above). Some volunteer opportunities include helping guide members of the public, posting and removing signs, and staffing the book sale table (if applicable). 

\subsubsection{Public Observing Nights and Events} 
The Anna McPherson observatory is located on the fifth floor of the Rutherford Building and hosts \textbf{monthly public observing nights} as well as other outreach and inreach events. The observatory coordinator is (historically) a graduate student volunteer role and the events are supported by undergraduate, graduate, postdoc, and staff volunteers who help operate telescopes and talk to the public about astronomy, current events in space, and astro-related research at McGill. Anyone in the department is welcome to volunteer for events, no prior experience required! Please email \href{mailto:observatory@physics.mcgill.ca}{observatory@physics.mcgill.ca} to be added to the volunteer list. The observatory coordinator monitors this inbox, but you can also (secondarily) reach out to the TSI inreach/outreach coordinator.  

The Outreach Committee often hosts outreach observation events during eclipses, such as providing educational materials and safe viewing methods for the 2026 partial, 2024 total solar eclipses, and rooftop observations of lunar eclipses (such as in 2025). These events are, naturally, seasonal and community involvement varies depending on the magnitude of the eclipse. For instance, the 2024 total solar eclipse was effectively a mass mobilization effort requiring total committee involvement, multiple graduate and undergraduate volunteers, an information fair, and featured a full crowd of visitors occupying the entire McGill Campus. Therefore, these events may require a large or small group of volunteers depending on popular demand.

\subsubsection{McGill Physics Hackathon}\index{Coordinators!Hackathon}
The McGill Physics Hackathon started in 2016 and is an annual computer programming competition open to anyone (and especially attended by undergraduate, CÉGEP, and high school students). The hackathon is an extensive endeavour, and organization efforts are led by the \textbf{Hackathon Outreach Coordinator} and the hackathon team. The hackathon also counts on the assistance of a team of volunteers, and is a popular way to become involved with science education and project development. To learn more about the physics hackathon, see the \href{https://www.hackathon.physics.mcgill.ca/}{McGill Physics Hackathon website}.

\subsubsection{Space Explorers}\index{Coordinators!Space Explorers}
Space Explorers is an Outreach Committee initiative through which recruited and trained volunteers visit elementary school classrooms to deliver hands-on, inquiry-based learning activities meant to initiate children to fundamental physics concepts (such as density, gravity, craters, and friction). These activities are designed for children aged 9-10 (grades 3-4) but can extend to younger (grade 2) and older (grades 5-6) children in some cases. Through Space Explorers, volunteers gain experience in teaching to children and developing their science communication skills, some activities may be in French so it might be an opportunity to practice.

\subsubsection{Science in Space} \index{Coordinators!Science in Space}
“Science in Space (in collaboration with Dell Technologies/Girls who Game) is a 10-week astronomy program introducing elementary students (grades 5-8; ages 10-14) to physics, astronomy, and telescopes. Students learn about telescopes through hands-on lessons, and ultimately design and build telescopes in Minecraft with the guidance of graduate student mentors (who themselves identify as members of marginalized groups in STEM). Science in Space is organized by the Science in Space Outreach Coordinator (see contact above). The program is aimed at girls and students of marginalized genders in STEM to mitigate attrition from STEM”.

\subsubsection{Science Fairs, Open Houses, and Tours}
There are many opportunities to volunteer at city-wide science fairs through the Outreach Committee, such as Family Science Day in September, and 24 Heures de Science in May, and AstroFest. Volunteer duties vary but often include ``classic" outreach activities relating to building a solar system mobile, aluminium foil boats, designing an extremophile, straw rockets, craters, constellations, and more. The physics department also has working relationships with several schools. For example, there is often demand for judges for the Royal Vale School science fair in the spring, and there is demand for outreach volunteers to work with ``World of STEM" participants.

Sometimes physics department members organize open houses to invite high school groups to learn more about the department, tour the Rutherford Physics Collection Museum (see below). In such cases, physics department members often seek a small group of graduate students to give brief talks about their research. 

Lastly, some laboratories and groups conduct their own outreach opportunities (e.g. the RAD kits designed by the McGill Radio Lab). Ask around to learn more!

\subsubsection{The Ground State Podcast (MGAPS)}
The Ground State Podcast is a recent outreach initiative started by the MGAPS council. \textbf{Note} that the podcast is \textbf{NOT} an Outreach Committee initiative but it is coordinated by the MGAPS liaison within the Outreach Committee, who serves as VP Professional Development within the MGAPS council. The podcast interviews students about their work with the aim of giving more visibility to the research done within the physics department. Specifically, the podcast primarily interviews students about published work and their experiences. If you wish to participate in the podcast initiative or be interviewed, reach out to the MGAPS liaison within the Outreach Committee.


\subsection{Departmental Divisions}
The McGill Physics Department is host to various research centers covering the primary research interests of the physics community. Many students within the physics department work primarily within one or multiple of these centers and their facilities. The MGAPS council is currently seeking to arrange a meeting of center organizers once a year to maintain an open communication channel between the centers.  

\subsection*{Center for High Energy Physics (CHEP)}\index{CHEP}
The Center for High Energy Physics (CHEP) is the primary center for people working in high energy and nuclear physics and is composed of subgroups in experimental high energy physics (XHEP), theoretical high energy physics (THEP), nuclear physics (NP), and particle astrophysics (PA). Founded in 1983, the CHEP is hosted in the \textbf{3rd level of the Rutherford Physics Building}, and features seminars and journal clubs covering its community's research interests including: Hadronic Physics Seminars, Particle and Astroparticle Seminars, Theory HEP Seminars, and HEP Theory Journal Club. 

The CHEP consists of faculty members affiliated with CERN, SNOLAB, TRIUMF, VERITAS, and RHIC, as well as a roster of \href{https://www.physics.mcgill.ca/chep/ras-f.html#XHEP}{research associates} and visiting scholars from various research centers and universities. The CHEP is also responsible for hosting the \href{https://www.physics.mcgill.ca/~francois/projects/guidelines.html}{Physics Computer Network}, which is utilized by various groups in the physics department including the CPM and the TSI. Researchers also frequently use clusters at the Digital Research Alliance of Canada (DRAC; formerly Compute Canada) to run demanding programs.

The CHEP is organized primarily by its faculty members, usually through the lead of a director, and seminars are run by select post-docs and doctoral students. The XHEP branch is also involved in organizing outreach initiatives, in collaboration with the Outreach Committee, including LEGO models, spark chambers, bubble chambers, and tours, see \href{https://www.physics.mcgill.ca/xhep/en/outreach/}{here} to learn more. For more information on the CHEP, see the \href{https://www.physics.mcgill.ca/chep/}{CHEP page} and see contacts below for faculty members primarily involved with CHEP organization. See also the \href{https://www.physics.mcgill.ca/xhep/en/}{XHEP} page for more information on experimental research. \textbf{Note} that the CHEP is currently in a transitional state following the retirement of its previous director so the contacts below are \textit{de facto} organizers. 

\vspace{-5 pt}
 \begin{table}[h!]
    \centering
    \renewcommand{\arraystretch}{1.2}
    \arrayrulecolor{ocre}
    \begin{tabular}{|>{\columncolor{ocre!20}}lc>{\columncolor{ocre!20}}c|}
    \hline
    \rowcolor{ocre!60}
  {\color{white}\large\sffamily \textbf{Role}} & {\color{white}\large\sffamily \textbf{Name}}   & {\color{white}\large\sffamily \textbf{Contact}} \\

        Faculty Member & Thomas Brunner & thomas.brunner@mcgill.ca \\
        Faculty Member & Brigitte Vachon & brigitte.vachon@mcgill.ca \\
        
        \hline
\end{tabular}
\arrayrulecolor{black}
\end{table}



\subsection*{Center for the Physics of Materials (CPM)}\index{CPM}
The Center for the Physics of Materials (CPM) was funded in 1989 and is the primary center for people working in condensed matter and materials physics. The center is present in the \textbf{4th level of the Rutherford Physics Building} and is funded by the \href{https://rqmp.ca/}{Regroupement Québécois sur les Matériaux de Pointe (RQMP)}, the Faculty of Science, and McGill University. The research scope of the center is incredibly broad and features members and associates from various departments including physics, chemistry, biology, computer science, materials engineering, electrical engineering, and chemical engineering. The center also employs research staff and technicians, and is directly involved with the \href{https://www.mcgill.ca/microfab/}{McGill Nanotools Micro- and Nanofabrication Facility} in the B1 level as well as the \href{https://www.physics.mcgill.ca/preplab/}{Sample Preparation Laboratory} (ERB 205-206) in the 2nd level of the Rutherford Physics Building.

The CPM is associated with three primary research networks: the RQMP, \href{https://www.intriq.org/}{Institut Transdisciplinaire d'Information Quantique (INTRIQ)}, \href{https://www.mcgill.ca/miam/}{McGill Institute for Advanced Materials (MIAM)}, and the McGill Quantum Center (MQC) (which is currently being built). Through these networks, CPM students attend student conferences including the Advanced Materials Academy (AMA) summer school, the INTRIQ seasonal conferences, MaQTech conferences, and the MQC conference. Through these networks, CPM members are able to utilize facilities in the Université de Montréal, the 
Polytechnique Montréal, and the Université de Sherbrooke, and coordinate internships with industry and other universities both in and out of Canada. Researchers also frequently use  clusters at the Digital Research Alliance of Canada (DRAC; formerly Compute Canada) to run demanding programs. Some CPM members also maintain third party affiliations within their own groups, one example being \href{https://nanoacademic.com/}{Nanoacademic Technologies}, and these affiliations serve as means for students to become involved with industry or other aspects of research. Beyond these affiliations, some CPM groups also collaborate and co-supervise with CHEP, TSI, and the Medical Physics Unit on particular projects.

The CPM is organized by a scientific directorship and a student committee, which organize weekly CPM seminars, the annual CPM day, and numerous social events, such as student talks and barbecues. The weekly CPM seminars feature invited speakers presenting research relevant to the center, while CPM Day is an annual conference organized by the CPM directorship and student committee. The student committee also coordinates the student talks, where students and postdoctoral researchers present and share their work with the CPM community. Any CPM student can become a member of the CPM committee and is encouraged to do so to contribute to community building. Given the interdisciplinary nature of the CPM, students involved with organizing will occasionally have to work with different institutes and departments. Students can volunteer to help organize CPM events, and volunteering work may include helping organize conferences, set up talks, and help gather materials for social events. Below is a table with contacts for the CPM directorship and student committee; for more information on CPM members and research visit the \href{https://cpm.research.mcgill.ca}{CPM website}.

\vspace{-5 pt}
 \begin{table}[hbt!]
    \centering
    \renewcommand{\arraystretch}{1.2}
    \arrayrulecolor{ocre}
    \begin{tabular}{|>{\columncolor{ocre!20}}lc>{\columncolor{ocre!20}}c|}
    \hline
    \rowcolor{ocre!60}
  {\color{white}\large\sffamily \textbf{Role}} & {\color{white}\large\sffamily \textbf{Name}}   & {\color{white}\large\sffamily \textbf{Contact}} \\

        Co-Director & Tami Pereg-Barnea & tamar.pereg-barnea@mcgill.ca \\
        Co-Director & Bradley Siwick & bradley.siwick@mcgill.ca \\
        Student Committee Member & Julia Azzi & julia.azzi@mail.mcgill.ca \\
        Student Committee Member & Clément Fortin & clement.fortin@mail.mcgill.ca \\
        Student Committee Member & Matheus A. S. Pessôa & matheus.pessoa@mail.mcgill.ca\\
        Student Committee Member & Louis Rossignol & louis.rossignol@mail.mcgill.ca\\
        
        \hline
\end{tabular}
\arrayrulecolor{black}
\end{table} 

\subsection*{Trottier Space Institute (TSI)}\index{TSI}
The Trottier Space Institute (TSI), founded in 2015 as the McGill Space Institute, is the primary center for people working in astrophysics, cosmology, and astronomy. The institute is an interdisciplinary center that brings together a broad range of space-related research areas including members from the departments of physics, earth and planetary science, atmospheric and oceanic sciences, and astrobiology. The TSI is located \textbf{beside the Rutherford Physics Building at 3550 Rue University}, and is funded by Matrox, the Power Corporation of Canada, the Faculty of Science, and McGill University. 

The TSI is associated with research networks and groups including the Canadian Astronomical Society (CASCA) and the Centre for Research in Astrophysics of Québec (CRAQ), as well as groups with specific themes such as the Trottier Institute for Research on Exoplanets (IREx) or the Canadian Epoch of Reionization Science Working Group. In addition, the TSI features researchers from a wide range of collaborations including: CHIME, CHORD, ALBATROS, MIST, HIRAX, SPT, LOFAR, SKA, JWST, and OMM. High-performance computing (HPC) access is key to much TSI research. Many research groups have allocations on Digital Research Alliance of Canada (DRAC; formerly Compute Canada) clusters (e.g. Fir, Trillium, Narval). Others use the Canadian Advanced Network for Astronomical Research (CANFAR), and some help develop the Canadian Data-Intensive Astrophysics PLatform (CanDIAPL). Some third parties affiliated with specific TSI groups include the McGill Arctic Research Station (MARS), Polar Continental Shelf Program (PCSP), and the National Research Council (e.g. Dominion Radio Astrophysical Observatory). Some TSI groups also active maintain relationships with industry, one of such cases being \href{https://www.t0.technology/}{t0.technology}, which may offer internships to graduate students.

The TSI is led by the TSI Board, which is headed by a director and a roster of TSI internal members, McGill internal members, and external members. The institute's activities are organized by coordinating staff and students. Student coordinators of lunches with TSI seminar speakers serve one-year terms. The graduate student representative to the TSI Board serves a two-year term, and may occasionally be invited to Faculty of Science-wide meetings with leaders of other institutes and departments. The TSI organizes discussion groups catered to the research interests present within the institute including: TSI Lunch Talks (general), Cosmo-ph (cosmology), Astro-ph (all astro), Transient Discussion (radio transients, both slow and fast), Black Hole Lunch, Planet Lunch (exoplanets), Neutron Star Discussion, Random Papers Discussion, and SpaceCraft. Another popular event is the \href{https://astronomyontap.org/locations/montreal-qc-canada/}{Astronomy on Tap} Montréal chapter, which offers biweekly events consisting of two short talks by members of the astronomy community, each followed by rounds of astronomy trivia. Some volunteer opportunities relate to planning, helping speakers craft their talks, helping manage the crowd, and creating and hosting trivia. See below for a list of contacts for the TSI Board and coordinators, also see the \href{https://tsi.mcgill.ca/}{TSI website} for more information on TSI activities and members.

\vspace{-5 pt}
 \begin{table}[hbt!]
    \centering
    \renewcommand{\arraystretch}{1.2}
    \arrayrulecolor{ocre}
    \begin{tabular}{|>{\columncolor{ocre!20}}lc>{\columncolor{ocre!20}}c|}
    \hline
    \rowcolor{ocre!60}
  {\color{white}\large\sffamily \textbf{Role}} & {\color{white}\large\sffamily \textbf{Name}}   & {\color{white}\large\sffamily \textbf{Contact}} \\

        Director & Victoria Kaspi & victoria.kaspi@mcgill.ca \\
        Inreach/Outreach Coordinator & Eesha Das Gupta & eesha.dasgupta@mcgill.ca \\
        Program Administrator & Alexandre David-Uraz & alexandre.david-uraz@mcgill.ca\\
        Computing Fellow & Siddarth Mohite & siddharth.mohite@mcgill.ca\\
        
        
        \hline
\end{tabular}
\arrayrulecolor{black}
\end{table} 


\subsection{Teaching Support Union at McGill (TSUM)}\index{TSUM}
The Teaching Support Union at McGill (TSUM) is the oldest Teaching Assistant Union in the province of Québec, and it represents graduate teaching assistants, invigilators, and course-based academic casuals (i.e. graders, course tutors, undergraduate student assistants, graduate student assistants, and teaching fellows) at McGill. TSUM was originally founded in January 11th, 1993 as the Association of Graduate Students Employed at McGill (AGSEM) but its name was changed in 2026 to reflect its growing number of undergraduate and non-student members. The TSUM is headquartered at 207-3641 Rue University, and operates as a highly organized workers union. The TSUM's highest decision-making body is the General Assembly, and General Assemblies for Unit 1 (teaching assistants), Unit 2 (invigilators), and Unit 3 (academic casuals) make decisions related to their unit. The combined General Assembly elects executives, amends the constitution, and approves the budget. 

In between General Assemblies, the Delegates' Council meets monthly to mandate the Executive Committee and organize union activity. Bargaining Committees are formed one year in advance of the expiration of a unit's respective collective agreement. Becoming involved with TSUM is an excellent starting point for graduate students to become familiar with university working conditions and structure. 

To become involved with TSUM, graduate students may be elected as TSUM delegates in the MGAPS General Assembly, simultaneously becoming TSUM representatives within the MGAPS council, or become elected for positions within the TSUM Executive Committee or Bargaining Committee, though such positions usually require the experience of first becoming a delegate. The TSUM has a very close relationship with MGAPS and its council, and TSUM delegates frequently organize to improve graduate students' wages and working conditions in the physics department. To get in touch with the physics TSUM delegates see the contacts below, and to learn more about TSUM visit the \href{https://www.agsem.ca}{TSUM website}.
\index{TSUM!Delegate}

 \begin{table}[h!]
 \hspace*{-0.2cm}
    \centering
    \renewcommand{\arraystretch}{1.2}
    \arrayrulecolor{ocre}
    \begin{tabular}{|>{\columncolor{ocre!20}}lc>{\columncolor{ocre!20}}cc|}
    \hline
    \rowcolor{ocre!60}
  {\color{white}\large\sffamily \textbf{Role}} & {\color{white}\large\sffamily \textbf{Name}}   & 
  {\color{white}\large\sffamily \textbf{Office}} &
  {\color{white}\large\sffamily \textbf{Contact}} \\
  TSUM Delegate 1 & Numa Karolinski & ERB 325 & \href{mailto:physics1.delegate@tsum-stsem.ca}{\color{black}physics1.delegate@tsum-stsem.ca} \\
    TSUM Delegate 2 & Varun Muralidharan & ERB 349 & \href{mailto:physics2.delegate@tsum-stsem.ca}{\color{black}physics2.delegate@tsum-stsem.ca} \\
    TSUM Delegate 3 & Eric Friesen Waldner & ERB 404 & \href{mailto:physics3.delegate@tsum-stsem.ca}{\color{black}physics3.delegate@tsum-stsem.ca} \\
        \hline  
\end{tabular}
\end{table}

\subsection{Post-Graduate Student Society (PGSS)}
The Post-Graduate Student Society (PGSS) of McGill University is the umbrella organization that encompasses all graduate student and post-doctoral associations, including MGAPS. The PGSS is a very extensive representative body that aims to provide resources to deal with university and Montreal life, advocate for students, improve student life, and plan social activities. The PGSS body politic is comprised of legislative, executive, and judicial branches under a board of directors, and featuring committees of council covering a variety of affairs, to learn more about PGSS structure see \href{https://pgss.mcgill.ca/en/get-involved}{here}.

The PGSS is headquartered in the \href{http://thomsonhouse.ca/home}{Thomson House}, a limestone mansion built in 1932 that was purchased by McGill University in 1968, and all graduate and post-doctoral fellows are granted associate membership. Thomson House serves as a common social club for events and dining. It is possible for graduate students and graduate student associations to reserve full rooms in Thomson House for events; however, it can come at a heavy fee and there are other options near campus that may provide better offers (such as the \href{https://bmp.coop/}{Co-Op Bar Milton-Parc}).

To become involved with PGSS, graduate students can become elected as PGSS representatives in the MGAPS council to represent the interests of physics graduate students in PGSS meetings and assemblies. Graduate students can also become directly involved with PGSS as volunteers, or become members of specific branches within the PGSS branches by applying for the positions or winning elections. These processes differ depending on the positions in question. To contact the PGSS physics representatives see the contacts below, and to learn more about PGSS visit the \href{https://pgss.mcgill.ca/en/}{PGSS Website}.

 \begin{table}[h!]
         \hspace*{-1cm}
    \centering
    \renewcommand{\arraystretch}{1.2}
    \arrayrulecolor{ocre}
    \begin{tabular}{|>{\columncolor{ocre!20}}lc>{\columncolor{ocre!20}}c|}
    \hline
    \rowcolor{ocre!60}
  {\color{white}\large\sffamily \textbf{Role}} & {\color{white}\large\sffamily \textbf{Name}}   & 
  {\color{white}\large\sffamily \textbf{Contact}} \\
  PGSS Representative 1 & Angela Sabzevari Gonzalez & \href{mailto:angela.sabzevarigonzalez@mail.mcgill.ca}{\color{black}angela.sabzevarigonzalez@mail.mcgill.ca} \\
    PGSS Representative 2 & Talia Martz-Oberlander & \href{mailto:talia.martz-oberlander@mail.mcgill.ca}{\color{black}talia.martz-oberlander@mail.mcgill.ca} \\
    PGSS Representative 3 & Valentin Boettcher & \href{mailto:valentin.boettcher@mail.mcgill.ca}{\color{black}valentin.boettcher@mail.mcgill.ca} \\
        \hline  
\end{tabular}
\end{table}

\subsection{Departmental Initiatives}
\subsubsection{Rutherford Physics Collection Museum}
Idealized by Prof. Jean Barrette, the Rutherford Physics Collection is a museum dedicated to displaying and maintaining historical physics instruments from the 19th and early 20th centuries, including the apparatus and documents from Ernest Rutherford’s time at McGill University. Recently, a paid position (with pay equivalent to a teaching assistant) became available to graduate students; however, the museum is currently in a transition period. Reinvigorating it may be an ideal chance to develop science communication and history of science skills. See the \href{https://www.physics.mcgill.ca/museum/}{website} for more information.

\subsubsection{McGill International Physicists’ Tournament (MIPT)}
The McGill International Physicists’ Tournament (MIPT) team, started in 2022 by Matheus Pessôa, is the inaugural Canadian team to participate in the International Physicists' Tournament (IPT). Every year, the team organizes physics problems in preparation for the IPT, and those involved get the chance to develop problem solving and education skills. The team is composed of Bachelor's and Master's students and is organized by an executive team. See the \href{https://mipt.sus.mcgill.ca/}{MIPT website} for more details on how to get involved.


\subsection{Medical Physics Unit} \index{Medical Physics}
The Medical Physics Unit is a completely separate department from the Department of Physics, featuring its own graduate student association and EDI Committee, and it is hosted within the Department of Oncology. However, the Medical Physics Unit does not have proper doctoral program, so graduate students pursuing doctoral degrees officially fall within programs in physics (PhD) or biomedical engineering (PEng). Graduate students pursuing doctorate degrees in the field of medical physics will still work with groups in the Medical Physics Unit but because some of these students are pursuing doctorates in physics, these students fall within MGAPS membership. This creates a particular quirk where only medical physics students in doctorate programs associated with physics can officially act as MGAPS members and members of the physics community. Regardless of this distinction, medical physics students can still become involved with the Department of Physics, most commonly by being selected by the \href{https://www.mcgill.ca/medphys/students/student-council}{Medical Physics Student Council} to act as MGAPS liaisons. Medical physics students officially associated with the Department of Physics through the aforementioned quirk can act as physics students and occupy the same positions as a regular student within the Department of Physics. The MGAPS council maintains an open doors policy with the Medical Physics Student Council through its liaisons such that all MGAPS events are open to medical physics students regardless of degree, and conversely the Medical Physics Student Council allows physics students to participate in select events in the Medical Physics Unit. For more information on the Medical Physics Unit, see the \href{https://www.mcgill.ca/medphys/}{Medical Physics Unit website}. Finally, Prof. Peter Watson from the Medical Physics Unit also runs a medical physics newsletter called the Dummy Source, see \href{https://thedummysource.github.io/}{here} for more information and to sign up.

\section{Technical Information} 
\subsection{Printing in Rutherford}\index{Printing}
There are a few ways to print in the Rutherford building.
\begin{enumerate}
  \item \textbf{uPrint}. This is a McGill-wide printing and scanning service which allows you to print off of any uPrint printers on campus. Once you have UPrint set up on your laptop, you can simply print from your laptop and use your student ID to swipe at any of the UPrint printers. Some more info:
  \begin{itemize}
      \item Note that this option costs money! Each B\&W print is \$0.077 and each colour print is \$0.27.
      \item The current uPrint printer in Rutherford is located on the 3rd floor.
      \item For instructions on how to get uPrint on your laptop, please visit this website: \href{https://mcgill.service-now.com/itportal?id=kb_article&sysparm_article=KB0011017}{uPrint services}
      \item Another simple no-setup alternative is to email your documents from your McGill email account to \textbf{uprint.mono@mcgill.ca} (for black \& white jobs) or \textbf{uprint.colour@mcgill.ca} (for colour jobs). More info \href{https://www.mcgill.ca/libraries/using-libraries/scan-print-copy}{here} (applicable campus wide, not just at the libraries).
  \end{itemize}
  \item \textbf{Physics Computer Network}. A detailed list of printers in Rutherford can be found \href{https://www.physics.mcgill.ca/~francois/projects/guidelines.html}{here}. You can also find \href{https://www.physics.mcgill.ca/~francois/projects/guidelines.html#printers}{instructions} on how to add and print using our Rutherford printers (note: it is FREE for you but your supervisor is charged for your usage so be mindful! The handbook committee suspects this isn't tracked directly unless you have submitted your credentials to be whitelisted on the physics network, a prerequisite to print in colour). For any question regarding setting up and using these printers, please contact the Computer System Operator (see \hyperref[contacts]{\textcolor{purble}{\ref*{contacts}}}). 
\end{enumerate} 

\subsubsection{Printing conference posters}
Although you can print your conference poster at basically any print shop in Montr\'eal, the easiest, closest, and relatively cheap way to do it is at CopiEUS. Some more info:
\begin{itemize}
    \item Location: McConnell Engineering Building, 1st floor
    \item Cost: \href{https://www.copieus.mcgilleus.ca/posters}{https://www.copieus.mcgilleus.ca/posters}
    \item More information about copiEUS (can do more than just print posters!): \href{https://www.copieus.mcgilleus.ca/}{copiEUS}
\end{itemize}

\subsection{Makerspace} \index{Makerspace}
There is a Makerspace located in the Wong building (Room 0080, basement level) that is open to support all your 3D printing and prototyping needs. The facility is equipped with: \index{3D printing}

\begin{itemize}
    \item 5 Prusa 3D printers (4 MK3 and 1 MK4 series)
    \item Other equipment that requires separate training with the technician:
    \begin{itemize}
        \item Formlabs SLA printer (Form 2) 
        \item Engraving machines (suited to "machine" prototype circuit boards)
    \end{itemize}
    \item Test equipment
    \begin{itemize}
        \item 3 Soldering stations \index{Soldering}
        \item Signal generator
        \item Power supplies, etc.
    \end{itemize}
\end{itemize}
    
All members of the department are welcome to use these resources, especially the 3D printers, once training is completed. Importantly, most of these resources are available to use for personal projects, not just for research.

\textbf{For New Users:}

To get started, you will need to sign up for a training session (approximately 1 hour) with the Makerspace TA. You can book a session using \href{https://bookings.cloud.microsoft/book/MakerspaceTraining1@McGill.onmicrosoft.com/?ismsaljsauthenabled}{this link}. After completing the training and signing the user agreement, you'll receive a card access to the facility, Monday to Friday, 9:00 AM - 5:00 PM.

\textbf{For Existing Users:}

If you'd like a refresher, feel free to use the link above to schedule another session.

For more details, visit the Makerspace \href{https://makerspace.physics.mcgill.ca/index.html}{website} to explore available resources and learn more about the Makerspace.

If you have any questions or need support with your 3D printing projects, don't hesitate to reach out!


    \subsection{Laptop Rentals} \index{Laptop Rentals}

The Computer Taskforce in the basement of Burnside (room 1B19) rents laptops out to students for 1 week at a time (free of charge), with possibility to renew as long as needed. Laptop rentals are requested by email to ctf.science@mail.mcgill.ca, more information about the rental process and what devices are available is provided \href{https://ctf.mcgill.ca/for_students/laptop_lending.html}{here}.
\\\\
You can also rent laptops and microphones through the McGill IT department. \href{https://mcgill.service-now.com/itportal?id=kb_article_view&sysparm_article=KB0010873#How}{This} webpage describes the equipment loan process and what you can borrow (free of charge).  Laptops are available for students for one month periods at most, and must be requested at least 2 business days in advance through the IT portal \href{https://mcgill.service-now.com/itportal?id=sc_cat_item&sys_id=6b72b78d133437085a165d122244b035&sysparm_category=6aee4c914f4a9200b18f59dd0210c75e}{here}. More information about laptop rentals in particular is provided on \href{https://mcgill.service-now.com/itportal?id=kb_article_view&sysparm_article=KB0010885}{this} webpage.

    \section{Mailing Lists and Chats}\index{Mailing Lists}

The Department of Physics and its student groups issue announcements and organize events through a variety of mailing lists and group chats. As a graduate student, you should already be added to the department-wide and general graduate physics student mailing lists but you must add yourself to specific group chats or mailing lists to receive further information about activities and seminars, as well as socialize with your peers. Below is a comprehensive list of these lines of communications and their purpose.

\subsection*{Mailing lists}

\noindent\textit{* = only to be used by the administration team}

\begin{itemize}
    \item Department-wide communications*: 
    \href{mailto:all-dept@physics.mcgill.ca}{all-dept@physics.mcgill.ca}

    \item Graduate physics student-wide communications: 
    \href{mailto:PhysicsGrads@campus.mcgill.ca}{PhysicsGrads@campus.mcgill.ca}

    \item Undergraduate physics student-wide communications: 
    \href{mailto:Physics_ugrads@campus.mcgill.ca}{Physics\_ugrads@campus.mcgill.ca}

    \item \href{https://www.physics.mcgill.ca/seminars/sem_lists.html}{Seminar notifications sign-up form}
\end{itemize}

\subsection*{Physics group chats}

\begin{itemize}
    \item \href{https://mcgillphysics-rag1431.slack.com/join/shared_invite/zt-3vlrg6luf-k0KqSG3LdCQ2V2P96XF_kQ}{Graduates and Post-Docs Slack}

    \item \href{https://join.slack.com/t/physoutreachv-8w11821/shared_invite/zt-46ado9t80-6DSn~Wb26pX5CK_io4bFZQ}{Outreach Slack}

    \item \href{https://discord.gg/6CQgQrdNrZ}{Graduates Sports Discord}

    \item \href{https://discord.gg/BEXgQhSnHJ}{Undergraduate Discord}

    \item \href{https://join.slack.com/t/mcgillphysics-jns2249/shared_invite/zt-3xqtkfyt0-kcsvyuS8yKqqAEUud70Vsg}{Department-wide summer activities Slack}
\end{itemize}

    